\paragraph{La normalisation d'alpha et sa configuration signée.}
\label{inc-ampl-clausura}
La coordonnée de $\alpha$ est déterminée dans la voie positionnelle par la composition des blocs régionaux et du registre dodécaphasique. Avant la division positionnelle, la combinaison orientée a pour composantes $Z_j=P_j+E_j-\Phi_j-K_j^{\mathrm{per}}$ et le bilan avec quotient entrant est $s_j=Z_j+c_{j+1}$. La normalisation satisfait
\begin{equation}
 A_j=s_j-1000c_j=Z_j+c_{j+1}-1000c_j,\qquad 0\leq A_j<1000.
 \label{inc-ampl-normalizacion}
\end{equation}
Cette identité conserve simultanément le bloc normalisé et la
contribution du quotient transporté. La limite des préfixes
compatibles détermine la coordonnée scalaire de la clôture. La
configuration signée conserve, quant à elle, les positions de la
combinaison antérieure à la normalisation qui présentent un défaut négatif.
Dans la fenêtre initiale, cette configuration est le vecteur $Z_0$ de
\eqref{exc:precarry}, et son support négatif est $H^-_\alpha$.

Deux contenus d'une même opération sont ainsi définis. La
normalisation positionnelle détermine une coordonnée réelle au moyen de ses
restes et quotients compatibles ; le support négatif détermine une
configuration de positions dans le cycle de douze composantes. Le signe
se rapporte au bilan des registres antérieur à la normalisation. Sa lecture
incidencielle reste disponible parce que l'état conserve ce bilan
avec le résultat normalisé. Cette articulation correspond à la
voie positionnelle démontrée ; la comparaison avec les opérateurs analytiques
de lacunes conserve la portée précise établie dans le chapitre
consacré à $\alpha$.
